\chapter{Lawvere's fixed-point theorem}

\section{Surjections and fixed points in categories}

\begin{definition}[Global element]
    Let $\mathcal C$ be a category with terminal object $1$. Given an object $A$ of $\mathcal C$, a \emph{global element} of $A$,
    or just \emph{element} of $A$, is a morphism $1 \to A$.
\end{definition}

\begin{definition}[Endomorphism]
    Let $\mathcal C$ be a category and $A$ an object of $\mathcal C$. An \emph{endomorphism} of $A$ is a morphism $f: A \to A$.
\end{definition}

\begin{definition}[Fixed point]
    In a category $\mathcal C$ with terminal object $1$, given an endomorphism $f: A \to A$, we say that $f$ has a \emph{fixed point} if there
    exists a global element $x: 1 \to A$ such that $f \circ x = x$.
\end{definition}

\begin{definition}[Surjection]
    In a category $\mathcal C$, given a morphism $f: A \to B$, we say that $f$ is a \emph{surjection} if there exists a morphism $s : B \to A$ such that $f \circ s = 1_B$.
\end{definition}

\begin{definition}[Point-surjective morphism]
    In a category $\mathcal C$ with terminal object $1$, given a morphism $f: A \to B$, we say that $f$ is \emph{point-surjective} if for every
    element $b: 1 \to B$ there exists an element $a: 1 \to A$ such that $f \circ a = b$.
\end{definition}

\begin{remark}[Surjective implies point-surjective]
    Obviously, in a category with terminal object, every surjection is point-surjective. This is trivial, since, given $f: A \to B$ a surjection and $b: 1 \to B$, we have
    \[ f \circ s = 1_B \implies f \circ (s \circ b) = 1_B \circ b = b, \]
    therefore we can take $s \circ b : 1 \to A$ as an element realizing $b$.
\end{remark}

\begin{definition}[Weakly point-surjective morphism]
    In a category $\mathcal C$ with terminal object $1$, given a morphism $f : A \to C^B$, if there exists the evaluation morphism $\ev_{C,B}$,
    we say that $f$ is \emph{weakly point-surjective} if for every element $g: 1 \to C^B$ there exists an element $a_g : 1 \to A$ such that, for every $b: 1 \to B$, we have
    \[ (f \circ a_g)^\sharp \circ b = g^\sharp \circ b. \]
\end{definition}

\begin{notation}
    We use the notation $(-)^\sharp$ to denote abbreviate the evaluation morphism, i.e.,\ if $h : 1 \to C^B$, then $h^\sharp := \ev_{C,B} \circ \langle h \circ !_B, 1_B \rangle : B \to C$.
    This way $h^\sharp \circ b = \ev_{C,B} \circ \langle h,b \rangle$ for every $b: 1 \to B$.
\end{notation}

\begin{remark}[Point-surjective implies weakly point-surjective]
    Let $\mathcal C$ be a category with terminal object $1$, given $f : A \to C^B$, if $f$ is point-surjective, then it is also weakly point-surjective.
    Indeed, given $g: 1 \to C^B$, since $f$ is point-surjective, there exists an element $a_g : 1 \to A$ such that $f \circ a_g = g$, therefore,
    applying $(-)^\sharp$ and precomposing with any $b$, we have the pointwise equality $(f \circ a_g)^\sharp \circ b = g^\sharp \circ b$.
\end{remark}

\begin{definition}[Diagonal morphism]
    In a category $\mathcal C$, given an object $A$, we call $\Delta_A : A \to A \times A$ the \emph{diagonal morphism} of $A$,
    defined as $\Delta_A := \langle 1_A, 1_A \rangle$.
\end{definition}

[AGGIUNGERE QUALCHE CONTROESEMPIO DI COSE CHE NON SONO ALTRE COSE]

\section{The fixed-point theorem}

\begin{theorem}[Lawvere's fixed-point]
    Let $\mathcal C$ be a cartesian closed category, if there exists a weakly point-surjective morphism $\phi : A \to B^A$, then every endomorphism $f: B \to B$ has a fixed point.
\end{theorem}

\begin{proof}
    Let $f : B \to B$ be an endomorphism.
    By the equivalence $\mathsf{CCC} \simeq \lambda\text{-}\mathsf{Calc}_0$, Theorem~\ref{thm:curry-howard-lambek-without-wnno-and-iterator}, we apply $\mathbf L : \mathsf{CCC} \to \lambda\text{-}\mathsf{Calc}_0$ to $\mathcal C$ and work in its internal language $\mathbf L(\mathcal C)$.
    By the definition of $\mathbf L$, the objects of $\mathcal C$ become types and its global elements become closed terms, composition becomes substitution, while evaluation and currying become application and abstraction.
    Using the fact that $\ulcorner f \urcorner : 1 \to B^B$ and $\ulcorner \phi \urcorner : 1 \to (B^A)^A$, and taking the respective closed terms in $\mathbf L(\mathcal C)$, we define the following closed term
    \[ q := \lambda_{x \in A} \ulcorner f \urcorner^\sharp \mathopen{\scalebox{1}[1.35]{$($}}\underbrace{\mathopen{\scalebox{1}[1.15]{$($}}\underbrace{\ulcorner \phi \urcorner^\sharp(x)}_{\makebox[0pt][c]{$\scriptstyle\text{function coded by }x$}}\mathclose{\scalebox{1}[1.15]{$)$}}^\sharp(x)}_{\text{applied to its own code}}\mathclose{\scalebox{1}[1.35]{$)$}} \in B^A. \]
    This closed term in $\mathbf L(\mathcal C)$ corresponds to the global element
    \[ q = \ulcorner f \circ \ev_{B,A} \circ (\phi \times 1_A) \circ \Delta_A \urcorner : 1 \to B^A, \]
    in $\mathcal C$. The weak point-surjectivity of $\phi$, expressed in $\mathbf L(\mathcal C)$, gives a closed term $p \in A$, equivalently an element $p : 1 \to A$, that can be regarded as a pointwise code of $q$, such that
    \[ \left(\ulcorner \phi \urcorner^\sharp(p)\right)^\sharp(a) \underset{\varnothing}{=} q^\sharp(a), \]
    for every closed term $a \in A$. Taking $a=p$, we define the closed term, $q$ applied to its own code
    \[ s := \left(\ulcorner \phi \urcorner^\sharp(p)\right)^\sharp(p) \underset{\varnothing}{=} q^\sharp(p) \in B. \]
    At the same time, using the above definition of $q$, we have
    \begin{align*}
        q^\sharp(p) &\underset{\varnothing}{=} (\lambda_{x \in A}\ulcorner f \urcorner^\sharp ((\ulcorner \phi \urcorner^\sharp(x))^\sharp(x)))^\sharp(p) \\
                    &\underset{\varnothing}{=} \ulcorner f \urcorner^\sharp((\ulcorner \phi \urcorner^\sharp(p))^\sharp(p)) &&(\text{$\beta$-reduction}) \\
                    &\underset{\varnothing}{=} \ulcorner f \urcorner^\sharp(s).
    \end{align*}
    Thus $\ulcorner f \urcorner^\sharp(s) \underset{\varnothing}{=} s$ in $\mathbf L(\mathcal C)$. We now apply $\mathbf C : \lambda\text{-}\mathsf{Calc}_0 \to \mathsf{CCC}$ to $\mathbf L(\mathcal C)$. By the definition of $\mathbf C$, closed terms determine arrows as equivalence classes of pairs, equality of terms gives equality of the corresponding classes, and composition is defined by substitution.
    Therefore, if $u \in 1$ is a variable, we find the following arrow equation in $\mathbf C(\mathbf L(\mathcal C))$:
    \[ [(x \in B,\ulcorner f \urcorner^\sharp(x))] \circ [(u \in 1,s)] = [(u \in 1,s)] \quad\text{in } \mathbf C(\mathbf L(\mathcal C)). \]
    Finally, the equivalence of categories $\mathsf{CCC} \simeq \lambda\text{-}\mathsf{Calc}_0$ provides us the natural isomorphism $\varepsilon : \mathbf C \circ \mathbf L \overset{\cong}{\Rightarrow }1_{\mathsf{CCC}}$,
    whose component at $\mathcal C$ is the isomorphism
    \[ \varepsilon_{\mathcal C} : \mathbf C(\mathbf L(\mathcal C)) \xrightarrow{\cong} \mathcal C. \]
    By the definition of $\varepsilon_{\mathcal C}$ given in the proof of \Cref{thm:curry-howard-lambek-without-wnno-and-iterator}, we have
    \[ \varepsilon_{\mathcal C}([(x \in B,\ulcorner f \urcorner^\sharp(x))])=f, \qquad \varepsilon_{\mathcal C}([(u \in 1,s)])=s, \]
    because $\ulcorner f \urcorner^\sharp(x) \underset{x}{=} \ev_{B,B} \circ \langle \ulcorner f \urcorner, x \rangle \underset{x}{=} \ev_{B,B} \circ \langle \ulcorner f \urcorner \circ !_B, 1_B \rangle \circ x \underset{x}{=} f \circ x$, using the definition of $(-)^\sharp$
    and the fact that $\mathcal C[x]$ is a CCC, so $\ulcorner f \urcorner^\sharp = f$; and $s \circ u \underset{u}{=} s \circ 1_1 = s$ in $\mathcal C[u]$, since $s$ is a closed term.
    Therefore, the functoriality of $\varepsilon_{\mathcal C}$ transports the preceding arrow equation to $f \circ s=s$ in $\mathcal C$, so $s$ is a fixed point of $f$.
\end{proof}

\section{Many fixed-points, one proof}
Let's now see how Lawvere's fixed-point theorem can be used to prove many other fixed-point theorems, diagonal arguments and even other useful results in logic and mathematics in general.
For each result, we will introduce the category needed to apply Lawvere's theorem, and we will show how the hypotheses of the theorem are satisfied in that category.

\subsection*{Cantor's theorem}
In $\mathsf{Set}$ we can prove Cantor's theorem, on the non-existence of surjections from a set to its power set.

\begin{theorem}[Cantor's theorem]
    Let $A$ be a set, then there is no surjection $f: A \twoheadrightarrow \mathcal P(A)$.
\end{theorem}

\begin{proof}
    Take $A \in \ob(\mathsf{Set})$, given the bijection $|2^A| = |\mathcal P(A)|$, we can equivalently prove that there is no surjection $f: A \twoheadrightarrow 2^A$.
    Let's assume, by contradiction, that there exists a surjection $f: A \twoheadrightarrow 2^A$. Since $\mathsf{Set}$ is cartesian closed,
    and surjections are weak point-surjective maps, by Lawvere's fixed-point theorem, every endomorphism $g : 2 \to 2$ has a fixed point.
    That is a contradiction, since the endomorphism $\neg : 2 \to 2$, defined as $\neg(0) = 1$ and $\neg(1) = 0$, has no fixed points.
\end{proof}

\subsection*{Cantor's diagonal argument}
By Cantor's diagonal argument, we mean the fact that $|\NN| < |\RR|$, proved through the Cantor's diagonalization.
Even though this result can be obtained as a corollary of Cantor's theorem, we can also prove it directly using Lawvere's fixed-point theorem.

\begin{theorem}[Cantor's diagonal argument]
    $|\NN| < |\RR|$.
\end{theorem}

\begin{proof}
    Let's consider the inclusion $\iota : 2^\NN \hookrightarrow \RR$, defined as
    \[ \iota : (a_i)_{i \in \NN} \mapsto \sum_{i \in \NN} a_i 10^{-(i+1)}. \]
    This is indeed injective, since if $(a_i)_{i \in \NN} \ne (b_i)_{i \in \NN}$, then there exists the minimum $i_0 \in \NN$ such that, WLOG, $a_{i_0} = 1$ and $b_{i_0} = 0$, therefore
    it easy to see that $\iota((a_i)_{i \in \NN}) - \iota((b_i)_{i \in \NN}) \geq \frac 89 10^{i_0+1} > 0$.\newline
    Let $s : \RR \to 2^\NN$ be a left-inverse of $\iota$, defined as
    \[ s : \RR \twoheadrightarrow 2^\NN : x \mapsto \begin{cases}
        \iota^{-1}(x) & \text{if $x \in \Img(\iota)$} \\
        (0)_{i \in \NN} & \text{otherwise}
    \end{cases}. \]
    Given an arbitrary map $f : \NN \to \RR$, define $\phi:=s \circ f : \NN \to 2^\NN$.
    We specialize the construction in the proof of Lawvere's fixed-point theorem to $\mathcal C=\mathsf{Set}$, $A=\NN$, $B=2$, the map $\phi$, and the endomorphism $\neg : 2 \to 2$.
    We construct the term
    \[ d:=\lambda_{i \in \NN}\ulcorner \neg \urcorner^\sharp((\ulcorner \phi \urcorner^\sharp(i))^\sharp(i)) \in 2^\NN, \]
    and, under the identification of the elements of $2^\NN$ with binary sequences, we have, by $\beta$-reduction, that
    \[ d^\sharp(i) \underset{\{i\}}{=} \ulcorner \neg \urcorner^\sharp((\ulcorner s \urcorner^\sharp(\ulcorner f \urcorner^\sharp(i)))^\sharp(i)). \]
    Thus $d$ is precisely the diagonal sequence: its $i$-th coordinate is the negation of the $i$-th coordinate of the $i$-th sequence $\ulcorner \phi \urcorner^\sharp(i)$.\newline
    If $d \in \Img(\phi)$, there would exist $p \in \NN$ such that $\ulcorner \phi \urcorner^\sharp(p) \underset{\varnothing}{=}d$.
    The calculation obtained by taking $a=p$ in the proof of Lawvere's theorem would then give
    \[ 2 \ni b:=(\ulcorner \phi \urcorner^\sharp(p))^\sharp(p) \underset{\varnothing}{=}d^\sharp(p) \underset{\varnothing}{=}\ulcorner \neg \urcorner^\sharp((\ulcorner \phi \urcorner^\sharp(p))^\sharp(p)) \underset{\varnothing}{=}\ulcorner \neg \urcorner^\sharp(b), \]
    contradicting the fact that negation has no fixed point in $2$, hence $d \notin \Img(\phi)$.
    Finally, let $x_d:=\ulcorner \iota \urcorner^\sharp(d) \in \RR$ and let's check that $x_d \notin \Img(f)$. If $x_d \underset{\varnothing}{=}\ulcorner f \urcorner^\sharp(k)$ for some $k \in \NN$, then
    \[ \ulcorner \phi \urcorner^\sharp(k) \underset{\varnothing}{=}\ulcorner s \urcorner^\sharp(\ulcorner f \urcorner^\sharp(k)) \underset{\varnothing}{=}\ulcorner s \urcorner^\sharp(x_d) \underset{\varnothing}{=}\ulcorner s \urcorner^\sharp(\ulcorner \iota \urcorner^\sharp(d)) \underset{\varnothing}{=}d, \]
    where the last equality holds since, in $\mathsf{Set}$, $s \circ \iota = 1_{2^\NN}$, so we have a contradiction.
    Therefore $x_d \notin \Img(f)$, and, since $f : \NN \to \RR$ was arbitrary, there is no surjection $\NN \twoheadrightarrow \RR$, because we constructed the diagonal element $x_d \in \RR$.
\end{proof}

\subsection*{Russell's paradox}

\section{Lawvere's fixed-point, cartesian form}
We can also prove a generalization of Lawvere's fixed-point theorem, that holds in a cartesian category, and it is equivalent to the original one in a cartesian closed category, because of the bijection given by [REFERENCE].

\begin{definition}
    Given a category $\mathcal C$ with terminal object $1$, we say that a morphism $f: A \times A \to B$ is \emph{weakly product point-surjective} if, for every morphism $g : A \to B$,
    there exists a morphism $a_g : 1 \to A$ such that $f \circ \langle a_g,a \rangle = g \circ a$ for every $a: 1 \to A$.
\end{definition}

\begin{theorem}[Lawvere's fixed-point, product form]
    Let $\mathcal C$ be a cartesian category, if there exists a weakly product point-surjective morphism $\phi : A \times A \to B$,
    then every endomorphism $f: B \to B$ has a fixed point.
\end{theorem}

\begin{proof}
    Let $f : B \to B$ be an endomorphism.
    By the equivalence $\mathsf{CC} \simeq \lambda\text{-}\mathsf{Calc}_\times$, \Cref{thm:curry-howard-lambek-cartesian-categories}, we apply $\mathbf L : \mathsf{CC} \to \lambda\text{-}\mathsf{Calc}_\times$ to $\mathcal C$ and work in its internal language $\mathbf L(\mathcal C)$.
    By the definition of $\mathbf L$, the objects of $\mathcal C$ become types, and the terms of type $B$ with free variables among $X$ are polynomial arrows $1 \to B$ in $\mathcal C[X]$; in particular, global elements become closed terms, and composition becomes substitution.
    Let $x : 1 \to A$ and $y : 1 \to B$ be indeterminates, therefore $x \in A$ and $y \in B$ are variables in $\mathbf L(\mathcal C)$. Let's define the term
    \[ q(x) := (f \circ y)[(\phi \circ \Delta_A \circ x)/y] \in B, \]
    with no free variables other than $x$. This term corresponds to the morphism
    \[ f \circ \phi \circ \Delta_A \circ x : 1 \to B \quad \text{in } \mathcal C[x]. \]
    Applying the weak product point-surjectivity of $\phi$ to $f \circ \phi \circ \Delta_A : A \to B$, gives a closed term $p \in A$, such that
    \[ \phi \circ \langle p,a \rangle \underset{\varnothing}{=} q(x)[a/x] \]
    for every closed term $a \in A$. Taking $a=p$, we define the closed term
    \[ s := \phi \circ \langle p,p \rangle \underset{\varnothing}{=} q(x)[p/x] \in B. \]
    Using the definition of $q(x)$ and substitution, we have
    \begin{align*}
        q(x)[p/x] &\underset{\varnothing}{=} (f \circ y)[(\phi \circ \Delta_A \circ x)/y][p/x] \\
                  &\underset{\varnothing}{=} (f \circ y)[(\phi \circ \Delta_A \circ p)/y] &&(\text{by substitution}) \\
                  &\underset{\varnothing}{=} (f \circ y)[s/y] \underset{\varnothing}{=} f \circ s
    \end{align*}
    Therefore, we have $f \circ s \underset{\varnothing}{=} (f \circ y) [s/y] \underset{\varnothing}{=} s$ in $\mathbf L(\mathcal C)$.
    We now apply $\mathbf C : \lambda\text{-}\mathsf{Calc}_\times \to \mathsf{CC}$ to $\mathbf L(\mathcal C)$. By the definition of $\mathbf C$,
    the substitution of terms corresponds to the composition of classes of arrows, as follows, for a variable $u \in 1$:
    \[ [(y \in B,f \circ y)] \circ [(u \in 1,s)] = [(u \in 1,s)] \quad\text{in } \mathbf C(\mathbf L(\mathcal C)). \]
    Finally, the equivalence of categories $\mathsf{CC} \simeq \lambda\text{-}\mathsf{Calc}_\times$ provides us the natural isomorphism $\varepsilon : \mathbf C \circ \mathbf L \overset{\cong}{\Rightarrow} 1_{\mathsf{CC}}$,
    whose component at $\mathcal C$ is the isomorphism
    \[ \varepsilon_{\mathcal C} : \mathbf C(\mathbf L(\mathcal C)) \xrightarrow{\cong} \mathcal C. \]
    By the definition of $\varepsilon_{\mathcal C}$, \Cref{thm:curry-howard-lambek-cartesian-categories}, we have
    \[ \varepsilon_{\mathcal C}([(y \in B,f \circ y)])=f, \qquad \varepsilon_{\mathcal C}([(u \in 1,s)])=s, \]
    since $f \circ y \underset{y}{=} f \circ y$ and $s \circ u \underset{u}{=} s \circ 1_1=s$ in the respective polynomial categories.
    Therefore, applying the functor $\varepsilon_{\mathcal C}$ to the preceding arrow equation gives $f \circ s=s$ in $\mathcal C$, so $s$ is a fixed point of $f$.
\end{proof}

\section{Other fixed-point theorems and diagonal arguments}
